% =============================================================================
% Section 8: Lean formalization and certificate verification
% Crisp distinction between Lean theorems / SAT certificates / audit / conditional.
% =============================================================================

\section{Lean formalization and certificate verification}
\label{sec:lean}

A central design decision of this work is to separate, with explicit
care, four distinct categories of mathematical claim:

\begin{enumerate}[leftmargin=*, label=(\textbf{C\arabic*})]
\item \textbf{Lean-certified theorems}: theorems whose proofs reduce
  to checkable Lean kernel terms in Mathlib, with $0$ \texttt{sorry}
  and $0$ project-introduced \texttt{axiom}.
\item \textbf{Independently verified computational certificates}:
  SAT/UNSAT or numerical certificates produced by an external tool
  (CaDiCaL, \texttt{cake\_lpr}) and importable into Lean via a
  trusted bridge.
\item \textbf{Empirical / audit evidence}: computational data
  (witness searches, kernel sentinels, sampled CSP runs) that informs
  the picture but does not itself constitute a formal proof.
\item \textbf{Conditional / conjectural discussion}: results stated
  under hypotheses (Hypothesis~H, GRH at $q = M$, the Hensel
  load-balancing conjecture) that we do not assume or prove.
\end{enumerate}

This section maps each named theorem of \cref{sec:main} to one of
these four categories, and describes the certificate-verification
pipeline.

\subsection{Formalized residue infrastructure}
\label{sec:lean-residue}

The residue-set infrastructure underlying
\cref{thm:stage1-reduction} is implemented in
\texttt{lean/Erdos647ResiduePartitionStage1.lean} and
\texttt{lean/Bridge2.lean}, which together provide:

\begin{itemize}[leftmargin=*]
\item the predicate \texttt{Bridge2.IsErdos647 : $\N \to$ Prop}
  encoding the $\tauf$-barrier condition;
\item the Stage-1 sieve as a Boolean predicate
  \texttt{Bridge2.survives : $\N \to$ Bool} on residues mod $M$;
\item the lemma
  \texttt{Bridge2.IsErdos647 $\Rightarrow$ survives ((2520 N) \% M)}
  (Lean theorem, kernel-checked).
\end{itemize}

The 96-element survivor set is enumerated by
\texttt{native\_decide}; the 41-element open subset is the residual
after subsequent eliminations carried out in
\texttt{Erdos647ResiduePartitionStage1} and the rootline modules.

\textbf{Category:} (C1) Lean-certified.

\subsection{The $v_{23}$ universal file}
\label{sec:lean-v23}

\Cref{thm:v23-partial-closure} is realized as 96 instances of the
\texttt{HenselRootline.hensel\_rootline\_closure} primitive in
\texttt{lean/generated/V23All96.lean}. Each instance uses the same
elementary divisor-count lemma
\texttt{tau\_1058\_mul\_prime\_gt4 : $\forall p, p$ prime $\to (1058 p)$.divisors.card $> 4$},
which is itself proved by \texttt{native\_decide} after a small case
analysis on $p \in \{2, 23\}$ vs.\ $p \notin \{2, 23\}$. The Mathlib
import surface is \texttt{Mathlib.NumberTheory.ArithmeticFunction}
plus \texttt{Mathlib.Data.Nat.Factorization.Basic}.

\textbf{Category:} (C1) Lean-certified.

\subsection{The 143-collapse verifier}
\label{sec:lean-143}

\Cref{thm:143-collapse} is fully Lean-formalized in
\texttt{lean/Erdos647\_143Collapse.lean}. The headline theorems are:

\begin{itemize}[leftmargin=*]
\item \texttt{openResidues41\_143\_collapse}: $\forall r \in
  \texttt{openResidues41},\, 143 \mid r$.
\item \texttt{openResidues41\_div\_11}, \texttt{openResidues41\_div\_13}:
  individual divisibility statements at each prime.
\item \texttt{U\_41\_explicit}: the explicit list
  $\Uset = \texttt{openResidues41.map (\textbackslash{}\textbackslash{} 143)}$
  in $\Z / 323\Z$, matching the 41-element subset stated in
  \cref{sec:q323-normalization}.
\item \texttt{openResidues41\_reconstruction}: the round-trip
  $143 \cdot u \bmod M = r$ for the canonical map $r \mapsto u$.
\end{itemize}

All proofs use \texttt{native\_decide}. Compile time on a standard
build was 187\,s; the \texttt{.olean} artifact is $\sim 122$\,KB.

\textbf{Category:} (C1) Lean-certified.

\subsection{The audit-chain SAT certificates}
\label{sec:lean-sat}

The finite-local insufficiency theorems
(\cref{thm:ple-insufficiency,thm:interval-cover-insufficiency}) rely
on SAT certificates rather than direct Lean proofs. The pipeline is:

\begin{enumerate}[leftmargin=*]
\item Encode the full safe-set CSP $\Phi_r(K_0, B, K_{\rm scan})$ in
  CNF (one Boolean variable per local valuation choice; clauses for
  the protected window, scan window, and safe-set constraints).
\item Solve with CaDiCaL; on SAT, emit the assignment.
\item Convert the run trace to LRAT format using \texttt{drat-trim}
  or \texttt{frat-rs}.
\item Verify the LRAT certificate with \texttt{cake\_lpr}, the
  CakeML-verified LRAT checker of \cite{TanHeuleMyreen2021}.
\item Import the verified result into Lean via
  \texttt{Std.Tactic.BVDecide}, which bit-blasts the encoded CNF
  to a Lean \texttt{BitVec} goal and discharges it through the
  same checker.
\end{enumerate}

The trust base is therefore: Lean kernel + Mathlib + the CakeML
\texttt{cake\_lpr} verifier. The original SAT solver (CaDiCaL) is
\emph{not} part of the trust base; it functions only as a hint
producer.

\textbf{Category:} (C2) Independently verified computational
certificate, importable as (C1).

\subsection{The structural-kill detector and CSP audit scripts}
\label{sec:lean-audit-scripts}

The structural-kill detector
(\texttt{scripts/interval\_cover\_structural\_kill.py}) and the CSP
audit chain
(\texttt{scripts/interval\_cover\_kernel\_space\_csp.py},
\texttt{scripts/interval\_cover\_full\_safe\_set\_csp.py},
\texttt{scripts/nested\_survivor\_tree\_csp.py}) are Python audit
scripts that produce the empirical evidence underpinning
\cref{rem:class-specific-struct-certs} and the audit summary
table of \cref{tab:audit-summary}. The independent re-derivation of
the small-shift kernel-class certificates
(\texttt{scripts/independent\_kill\_verifier.py}, with descending
prime order and a fresh implementation) confirms that the certificates
are deterministic facts about the JSON kernel data, not artifacts of
any single solver run.

\textbf{Category:} (C3) empirical / audit evidence. The
mathematical content (per-kernel-class kill, CSP escape) is
recorded as part of the audit and \emph{not} promoted to a
Lean-certified theorem in this paper; it is precisely the
``not base-AP'' qualifier of \cref{tab:audit-summary} that
distinguishes this category from (C1)/(C2).

\subsection{The conditional barrier theorem}
\label{sec:lean-conditional}

\Cref{thm:conditional-barrier} is stated under Schinzel's
Hypothesis~H. We do not assume H elsewhere in the paper, and the
theorem is not Lean-formalized as a \emph{conditional} theorem
because Mathlib does not currently include a formalization of
Hypothesis~H. The structural admissibility of the polynomial system
$\Sigma_u$ for each $u \in \Uset$ is independently checkable (and
is a (C1) Lean theorem in principle, though not currently stated as
such); the implication ``H $\Rightarrow$ infinite-survivor'' is
classical and folklore.

\textbf{Category:} (C4) Conditional, with the underlying admissibility
checks in (C1).

\subsection{Trust-chain summary}
\label{sec:lean-trust-chain}

\Cref{tab:trust-chain} summarizes the categorization of every named
theorem of \cref{sec:main}.

\begin{table}[h]
\centering
\begin{tabular}{lll}
\toprule
\textbf{Theorem} & \textbf{Category} & \textbf{Lean / certificate file} \\
\midrule
\cref{thm:stage1-reduction} (96 residues) & (C1) Lean & \texttt{Bridge2.lean, ResiduePartitionStage1.lean} \\
\cref{thm:v23-partial-closure} ($v_{23}$) & (C1) Lean & \texttt{V23All96.lean} (96 instances) \\
\cref{thm:143-collapse} (143-collapse) & (C1) Lean & \texttt{Erdos647\_143Collapse.lean} \\
\cref{thm:ple-insufficiency} (\PLEclass) & (C2) cert & SAT $\Phi_r$ + LRAT + \texttt{cake\_lpr} \\
\cref{thm:interval-cover-insufficiency} & (C2) cert & full safe-set $\Phi_r$ certs \\
\cref{rem:class-specific-struct-certs} & (C3) audit & detector / CSP audit scripts \\
\cref{thm:conditional-barrier} (Schinzel-H) & (C4) cond. & admissibility (C1); H not in Mathlib \\
\cref{prop:empirical-r0} (empirical) & (C3) audit & witness search infrastructure \\
\bottomrule
\end{tabular}
\caption{Categorization of named theorems by trust chain. (C1) = Lean
kernel + Mathlib; (C2) = (C1) plus an independent verifier
(\texttt{cake\_lpr}); (C3) = empirical / audit evidence; (C4) =
conditional on a stated hypothesis.}
\label{tab:trust-chain}
\end{table}

\subsection{What is intentionally not formalized}
\label{sec:lean-not-formalized}

Several pieces of mathematical content in this paper are
\emph{intentionally not Lean-formalized}, with reasons:

\begin{itemize}[leftmargin=*]
\item \textbf{Schinzel's Hypothesis H} (used in
\cref{thm:conditional-barrier}): not in Mathlib; formalizing the
hypothesis is a project in its own right.
\item \textbf{The Selberg upper-bound divergence lemma}
(\cref{lem:selberg-divergence}): elementary calculus on a stated
upper bound; pen-and-paper rigor is sufficient and the result does
not load-bear formal certification.
\item \textbf{The empirical floor at the principal residue}
(\cref{prop:empirical-r0}): the witness search verifier is a
separately documented Python program; formalization of an
exhaustive numerical search is not the appropriate locus of formal
proof for this paper.
\item \textbf{The Bateman--Horn singular series predictions} in
\cref{thm:conditional-barrier}: classical, conditional on H, and
not load-bearing for the barrier.
\end{itemize}

The boundary between (C1)/(C2) and (C3)/(C4) is not arbitrary: the
positive structural reductions
(\cref{thm:stage1-reduction,thm:v23-partial-closure,thm:143-collapse})
and the SAT-certificate-importable insufficiency theorems
(\cref{thm:ple-insufficiency,thm:interval-cover-insufficiency})
constitute the kernel-checked content of the paper; the empirical
floor and conditional discussion are presented as separate
categories without claiming kernel-level certification.

% =============================================================================
% End of section 8
% =============================================================================
